% ============================================================

% ============================================================

\documentclass[12pt, onecolumn]{IEEEtran}

\usepackage[T1]{fontenc}
\usepackage[utf8]{inputenc}
\usepackage{graphicx}
\usepackage{booktabs}
\usepackage{array}
\usepackage{tabularx}
\usepackage{xcolor}
\usepackage{hyperref}
\usepackage{cite}
\usepackage{caption}
\usepackage{enumitem}
\usepackage{microtype}
\usepackage{colortbl}
\usepackage{geometry}
\geometry{margin=2.5cm}

\definecolor{rowgray}{RGB}{242,242,242}
\definecolor{ieeblue}{RGB}{31,78,121}

\hypersetup{colorlinks=true, linkcolor=ieeblue,
            citecolor=ieeblue, urlcolor=ieeblue}
\captionsetup{font=small, labelfont=bf, textfont=it,
              skip=5pt, justification=centering}

% ============================================================
\begin{document}

% ── Title ────────────────────────────────────────────────────
\title{Community Security Alert System (CSAS):\\
       Systems Analysis and Design --- Workshop No.~1}

\author{
\IEEEauthorblockN{
Felipe Jose Garzon Herrera\_20251020132\\
Quintero Gordillo Juan Esteban --- 20251020137\\
Garrido Medina Henry Samuel\\
Cusba Mar\'{i}n Gabriel Mateo
}
\IEEEauthorblockA{
Computer Engineering Program\\
Universidad Distrital Francisco Jos\'{e} de Caldas\\
Bogot\'{a} D.C., Colombia --- Semester 2026-I
}
}

\maketitle

% ── Abstract ─────────────────────────────────────────────────
\begin{abstract}
This report presents a complete systems analysis of the Community
Security Alert System (CSAS), a crowd-sourced platform designed to
improve safety on university campuses. The document covers system scope
and boundary definition, stakeholder identification, operational
environment characterization, data collection planning and execution,
data quality evaluation, system element decomposition, relationship
analysis, and a sensitivity study of key behavioral parameters. Primary
data was collected through structured interviews with campus security
staff, student surveys ($n=20$), direct observation across four campus
zones (8 sessions total), document analysis, and benchmarking of
existing security platforms in Bogot\'{a}. Findings consistently show
that night-time incident concentration, a low community reporting
culture, and the absence of a centralized mobile reporting tool are the
main gaps the CSAS is designed to address.
\end{abstract}

\begin{IEEEkeywords}
campus security, community alert system, systems analysis,
crowd-sourced incident reporting, real-time notifications,
data collection, sensitivity analysis, feedback loops
\end{IEEEkeywords}

% ============================================================
%  I. SYSTEM OVERVIEW AND SCOPE DEFINITION
% ============================================================
\section{System Overview and Scope Definition}

\subsection{System Description}

The Community Security Alert System (CSAS) is a collaborative
technological platform designed to enhance safety within university
campuses and their immediate surroundings. By integrating real-time
notifications with crowd-sourced incident reporting, the system empowers
students, faculty, and staff to act as active sensors for community
safety. Users can report incidents---ranging from medical emergencies
and accidents to suspicious activities---via a mobile or web interface.
To maintain data integrity, the system employs verification mechanisms
including multi-user confirmation and administrative moderation,
ensuring that alerts are both high-speed and high-reliability.

\subsection{System Boundaries}

Defining clear boundaries is essential for scoping the analysis and
preventing scope overextension in subsequent design phases.
Table~\ref{tab:boundaries} separates what falls within the CSAS from
what constitutes its external environment.

\begin{table}[h]
\centering
\caption{System Boundary Definition --- CSAS}
\label{tab:boundaries}
\renewcommand{\arraystretch}{1.3}
\begin{tabular}{>{\bfseries\small}p{0.44\linewidth}
                >{\small}p{0.44\linewidth}}
\toprule
\textbf{Inside the system (scope)} &
\textbf{Outside the system (environment)} \\
\midrule
Incident reporting \& management platform
    & University physical infrastructure \\
\rowcolor{rowgray}
Crowd-sourced verification logic
    & Local law enforcement (Police / EMS) \\
Real-time push notification engine
    & Public telecommunications \& ISP networks \\
\rowcolor{rowgray}
User and incident databases
    & Non-registered external visitors \\
Emergency coordination dashboard
    & Campus security physical hardware (CCTV, etc.) \\
\bottomrule
\end{tabular}
\end{table}

\subsection{Stakeholder Identification}

Six stakeholder groups interact with the CSAS. \textbf{Students} are
the primary users and the main contributors to crowd-sourced incident
reporting; they submit reports of suspicious activities and receive
real-time zone-specific notifications. \textbf{Professors and
administrative staff} use the system to remain informed about campus
safety conditions, reporting situations that may threaten the
community. \textbf{The campus security team} plays the critical
operational role: they review and verify reports, assess severity,
coordinate the appropriate response, and monitor security trends across
campus. \textbf{Local authorities}---police departments and emergency
services---receive verified incident data when incidents require
external intervention. \textbf{System administrators} maintain the
technical operation of the platform, managing accounts, performance,
and security. \textbf{University management} oversees the strategic use
of the CSAS, consuming analytics and incident statistics to allocate
resources and guide safety policy decisions.

\subsection{System Purpose}

The main purpose of the CSAS is to enhance campus safety through
real-time communication and collaborative incident reporting. The
system aims to: enable crowd-sourced reporting of security incidents;
provide real-time safety notifications for specific campus areas;
improve awareness of potential security risks; facilitate verification
of reported incidents; support rapid emergency response coordination;
and encourage active community participation in campus safety.

\subsection{Operational Environment}

The CSAS operates across four interconnected environmental dimensions.
The \textbf{technological} dimension requires mobile and web clients,
internet connectivity, secure cloud servers, and a real-time
notification delivery infrastructure. The \textbf{social} dimension
encompasses a heterogeneous community of students, professors, staff,
and visitors who both generate and consume safety information. The
\textbf{physical} dimension covers the university campus and adjacent
neighborhoods, including academic buildings, parking lots, residences,
and public thoroughfares. The \textbf{institutional} dimension involves
coordination with university administration, the campus security
department, and local law-enforcement agencies.

\subsection{System Context Diagram}

Fig.~\ref{fig:context} presents the system context diagram, showing the
information flows between the CSAS and its external actors. Students
and professors submit incident reports and receive real-time alerts.
The campus security team performs monitoring, verification, and
emergency coordination, forwarding escalated cases to local authorities.
University management consumes security reports and analytics for
strategic decision-making.

\begin{figure}[h]
  \centering
  \includegraphics[width=0.72\linewidth]{context_diagram.png}
  \caption{System context diagram --- CSAS. Arrows indicate the
  direction of information flow between each stakeholder group and
  the platform. \textit{Source: own elaboration (2026).}}
  \label{fig:context}
\end{figure}

\subsection{System Architecture Diagram}


Fig.~\ref{fig:arch} presents the full system architecture as a node
relationship graph. The diagram organizes all components across four
groups: external actors (students, security staff, university admin,
police, and maintenance), the interface layer (mobile app, web portal,
SMS gateway, and push alerts), the processing core (REST API, AI
verifier, alert dispatcher, geolocation service, and auth module), and
the data layer (incidents database, cache, media store, analytics
engine, and audit logs). Edges between nodes represent the primary
information flows described throughout this document.

\begin{figure}[h]
  \centering
  \includegraphics[width=\linewidth]{arch_nodes.png}
  \caption{System architecture --- CSAS node relationship graph.
  Node size reflects relative centrality within the system; the REST
  API is the highest-connectivity node, acting as the central gateway
  between all layers. \textit{Source: own elaboration (2026).}}
  \label{fig:arch}
\end{figure}

% ============================================================
%  II. DATA COLLECTION PLAN AND EXECUTION
% ============================================================
\section{Data Collection Plan and Execution}

\subsection{Data Collection Plan}

To analyze the CSAS, we designed a data collection plan that combined
multiple methods simultaneously. The goal was not to rely on a single
source, but to bring together different perspectives to build a more
complete picture of the problem: what students think, how the security
staff works, and what actually happens on campus at different times of
day. Table~\ref{tab:plan} summarizes the six methods applied.

\begin{table}[h]
\centering
\caption{Data Collection Plan --- Six Methods Applied to the CSAS}
\label{tab:plan}
\renewcommand{\arraystretch}{1.3}
\begin{tabular}{>{\bfseries\small}p{0.14\linewidth}
                >{\small}p{0.16\linewidth}
                >{\small}p{0.16\linewidth}
                >{\small}p{0.10\linewidth}
                >{\small}p{0.30\linewidth}}
\toprule
\textbf{Method} & \textbf{Tool} & \textbf{Sample} &
\textbf{Duration} & \textbf{Objective} \\
\midrule
Interviews
  & Structured interview guide
  & 5 security staff, 2 administrative staff
  & 1 week
  & Understand how incident reporting currently works, identify
    the main problems experienced by security personnel, and
    determine what a new system would need to be useful to them. \\
\rowcolor{rowgray}
Surveys
  & Google Forms / Microsoft Forms
  & 20 students
  & 2 weeks
  & Assess how safe students feel on campus, whether they have
    experienced incidents, and whether they would use an alert
    application. \\
Direct observation
  & Observation record form
  & 4 campus zones
  & 1 week
  & Record what happens in different parts of the campus at
    different times of day: how many people are present, how
    active each zone is, and whether anything suspicious occurs. \\
\rowcolor{rowgray}
Document analysis
  & Institutional reports and security records
  & Existing campus incident reports
  & 3 days
  & Review official incident records to identify historical
    patterns and compare them with what we found in the field. \\
Benchmarking
  & Comparative analysis
  & 3 security apps in Bogot\'{a}
  & 1 week
  & Study what existing platforms do well and where they fall
    short, to avoid reinventing the wheel and improve upon
    what already exists. \\
\rowcolor{rowgray}
Weather monitoring
  & Weather apps
  & Teusaquillo locality
  & 1 week
  & Explore whether there is a relationship between weather
    conditions (rain, temperature, time of day) and the levels
    of activity or incidents observed on campus. \\
\bottomrule
\end{tabular}
\end{table}

\subsection{Interviews}

We spoke with five security staff members and two administrative staff
members. The objective was to understand the day-to-day reality of
campus security from the inside, not only from the student perspective.
The results were clear: staff identified three main problems that came
up in nearly every interview. First, communication times are slow ---
from the moment an incident occurs to when the response team is notified
can take several minutes, mainly because reporting depends on someone
calling or physically approaching a security guard. Second, there is no
centralized reporting tool; each incident is handled differently with no
systematic record that allows pattern analysis or follow-up. Third,
staff expressed a clear need for faster notifications that include the
exact location of the incident, not just a verbal description.

\subsection{Student Survey}

We distributed the survey to 20 students through an online form. The
questions measured three things: how safe students feel on campus,
whether they have experienced any incident, and whether they would use
an alert app. Table~\ref{tab:survey} and Fig.~\ref{fig:survey} present
the data collected and the main results.

\begin{table}[h]
\centering
\caption{Raw survey data (n=20). Variables: safety perception,
incident experience, incident type, willingness to use the app,
and time of day.}
\label{tab:survey}
\includegraphics[width=0.95\linewidth]{survey_table.png}
\end{table}

\begin{figure}[h]
  \centering
  \includegraphics[width=0.52\linewidth]{survey_chart.png}
  \caption{Proportion of students who experienced a security incident
  (n=20). 35\% reported having experienced or witnessed an incident.
  \textit{Source: own elaboration (2026).}}
  \label{fig:survey}
\end{figure}

The chart shows that 7 of the 20 students surveyed (35\%) reported
having experienced or witnessed some type of incident on or around
campus. Although the sample is small, this proportion is notably high
and is consistent with what security staff mentioned during interviews.
Three patterns stand out from the data. First, all reported incidents
(theft, harassment, and suspicious activity) occurred almost entirely
during afternoon and night hours, confirming that the evening period is
the most vulnerable. Second, 100\% of students who experienced an
incident said they would use the alert app, compared to 69\% among
those who had not experienced anything --- this makes sense, since
direct experience with an incident makes the utility of the system
immediately obvious. Third, several students who said they felt safe on
campus had also reported experiencing an incident, which may reflect
normalization of risk or incidents occurring outside campus on
commuting routes --- a question worth investigating further.

\subsection{Direct Observation}

We conducted observation sessions in four campus zones: the main
entrance, the library area, the parking lot, and the student residences.
Each zone was observed during at least two different time slots to
compare behavior according to the time of day. Table~\ref{tab:obs}
and the accompanying analysis summarize the results.

\begin{table}[h]
\centering
\caption{Direct observation record --- 8 sessions across 4 campus
zones. Variables: activity level and presence of suspicious event.}
\label{tab:obs}
\includegraphics[width=0.82\linewidth]{observation_table.png}
\end{table}

Of the eight sessions, only two recorded suspicious events and both
were in the parking lot at night. The other zones --- especially the
main entrance and the library --- showed high or medium activity levels
with no irregularities, likely because higher foot traffic and better
lighting provide an informal deterrent effect. Four key findings emerge
from this phase. The parking lot at night is clearly the highest-risk
observed zone, with poor lighting and low foot traffic as contributing
factors. The student residence area showed poor lighting at two specific
points, even though no suspicious event was recorded during observation;
this represents a latent risk that does not appear in official records.
Additionally, the suspicious events we observed were not reported by any
other witness during the observation period, which reinforces the
under-reporting hypothesis: incidents happen but people do not report
them.

\subsection{Document Analysis}

We reviewed the available institutional security reports and some
publicly accessible data on campus incidents. The goal was to look for
historical patterns that could confirm or contradict what we were
finding through the other methods. The most striking finding was that
formal records show far fewer incidents than actually occur according to
the survey and observation data. This confirms a significant
under-reporting problem: people do not report what they see. The most
documented incident types are theft and unauthorized access, which match
the types mentioned in the student survey.

\subsection{Benchmarking}

We reviewed three existing community alert platforms operating in
Bogot\'{a} to understand what works, what does not, and where there is
an opportunity to improve with the CSAS. Table~\ref{tab:bench} presents
the comparison.

\begin{table}[h]
\centering
\caption{Benchmarking of existing security alert platforms in Bogot\'{a}.}
\label{tab:bench}
\renewcommand{\arraystretch}{1.3}
\begin{tabular}{>{\bfseries\small}p{0.18\linewidth}
                >{\small}p{0.22\linewidth}
                >{\small}p{0.22\linewidth}
                >{\small}p{0.26\linewidth}}
\toprule
\textbf{Platform} & \textbf{Strengths} & \textbf{Weaknesses} &
\textbf{What we can take for CSAS} \\
\midrule
Policia Nacional / Vigila App
  & Direct integration with police; national coverage
  & Slow response; depends on user placing a phone call
  & The idea of direct authority integration is key. CSAS should have
    something similar but faster. \\
\rowcolor{rowgray}
Neighborhood networks (WhatsApp)
  & Easy to use; very high adoption; instant notifications
  & No verification, no records, prone to generating panic
  & Immediacy and channel familiarity are good models, but verification
    must be added to avoid false alarms. \\
SIURE UD (internal system)
  & Already exists; staff are familiar with it
  & Email and phone only; no mobile app or real-time capability
  & CSAS can complement it directly rather than replace it, building
    on the existing institutional process. \\
\bottomrule
\end{tabular}
\end{table}

\subsection{Weather Monitoring}

As an additional method, we tracked weather conditions in the Teusaquillo
locality for one week and tried to cross-reference that information with
the activity levels observed on campus. The hypothesis was simple: when
it rains at night, fewer people are outside and that may increase risk.
The results are preliminary since one week is not enough time to draw
definitive conclusions, but we did notice that both suspicious events
recorded during the observation phase coincided with rainy nights. It
will make sense to explore this correlation with more data in subsequent
phases of the project.

% ============================================================
%  III. DATA QUALITY EVALUATION
% ============================================================
\section{Data Quality Evaluation}

\subsection{Identified Limitations}

Although the data we collected gives a useful picture of the system, it
is important to be honest about its limitations. This does not mean the
work is flawed --- it means we know how far our conclusions reach and
where we need to go deeper. Table~\ref{tab:limits} summarizes the main
limitation identified for each method and what it implies for the
interpretation of results.

\begin{table}[h]
\centering
\caption{Limitations by data collection method and their implications.}
\label{tab:limits}
\renewcommand{\arraystretch}{1.3}
\begin{tabular}{>{\bfseries\small}p{0.15\linewidth}
                >{\small}p{0.26\linewidth}
                >{\small}p{0.45\linewidth}}
\toprule
\textbf{Method} & \textbf{Limitation} & \textbf{Implication} \\
\midrule
Survey
  & Sample of only 20 students
  & Results cannot be generalized to the whole university. They are
    a useful first approximation for identifying trends, but are not
    statistically representative. \\
\rowcolor{rowgray}
Survey
  & Self-reported responses (unverified)
  & Some students may have answered what they thought was expected
    rather than what they actually think or experienced. The anonymous
    format helps reduce this effect. \\
Observation
  & Only 4 zones and 8 sessions
  & Parts of campus were not observed: cafeterias, bathrooms, indoor
    corridors. Patterns from those areas were not captured. \\
\rowcolor{rowgray}
Observation
  & Observer effect
  & Some behaviors may have changed because someone was taking notes.
    We tried to be unobtrusive, but this is a limitation that is
    difficult to eliminate entirely. \\
Document analysis
  & Limited access to detailed data
  & Official reports do not include breakdowns by zone or time slot,
    so we could not perform as specific a historical analysis as we
    would have liked. \\
\rowcolor{rowgray}
Benchmarking
  & Public information only
  & We did not have access to real usage data or internal metrics
    from the compared platforms, only their official documentation. \\
Weather monitoring
  & Only one week of observation
  & Not enough to establish a real correlation between weather and
    security. Findings are indicative and will need validation with
    more data over a longer period. \\
\bottomrule
\end{tabular}
\end{table}

\subsection{Why the Data Remains Useful}

Despite these limitations, we consider the collected data sufficient
for the objectives of this workshop phase. Three reasons support this
assessment. First, we used multiple methods simultaneously: the most
important findings --- incidents concentrated at night, low reporting
levels, preference for a mobile channel --- appear independently in the
interviews, the survey, and the direct observation. When three different
methods converge on the same conclusion, there is stronger confidence
that it reflects reality. Second, all data is primary and collected
directly by the team, which gives us full control over the collection
process and allows us to trace the origin of every data point. Third,
the data is sufficient to guide the system design: we do not need a
sample of 500 people to know that the parking lot at night is a critical
risk point or that people do not know how to report an incident. That is
already clear from what we have. The identified limitations do not
invalidate the work; they simply define clearly how far our conclusions
reach and set the research agenda for the next validation phase.

% ============================================================
%  IV. ELEMENT AND RELATIONSHIP ANALYSIS
% ============================================================
\section{Element and Relationship Analysis}

\subsection{System Elements}

The CSAS is composed of five core elements that work together to
improve safety in the campus environment. \textbf{Users} (students and
staff) are the primary actors who report incidents and receive alerts;
they are the main data-input layer of the system. The \textbf{reporting
platform} is the mobile or web interface through which users submit
incident information including location, time, type, and description.
\textbf{Campus security staff} review incoming reports, verify their
content, and coordinate the appropriate response; interviews with
security personnel revealed that current systems suffer from
communication delays, which highlights the need for faster and more
centralized tools \cite{sierra2026}. The \textbf{notification system}
broadcasts verified alerts to other users in affected areas, enabling
proactive risk awareness. The \textbf{incident database} stores all
incidents, verifications, and outcomes, enabling pattern analysis and
historical reference over time.

\subsection{System Relationship Map}

Fig.~\ref{fig:relmap} illustrates how information flows between the
five core system elements. Users report incidents through the reporting
platform, which forwards each report to the verification process and to
campus security staff. After verification, alerts are dispatched through
the notification system back to nearby users. All incident data---from
initial submission to final resolution---is logged in the incident
database for future analysis.

\begin{figure}[h]
  \centering
  \includegraphics[width=0.60\linewidth]{relationship_map.png}
  \caption{System relationship map --- CSAS. Information flows from
  users through the reporting platform to the verification process and
  security staff; verified alerts are dispatched via the notification
  system while all data is logged in the incident database.
  \textit{Source: own elaboration (2026).}}
  \label{fig:relmap}
\end{figure}

% ============================================================
%  V. SENSITIVITY, COMPLEXITY, AND EMERGENT BEHAVIOR
% ============================================================
\section{Sensitivity, Complexity, and Emergent Behavior}

\subsection{Parameter Sensitivity Analysis}

The behavior of the system can change significantly depending on several
operational parameters. Table~\ref{tab:sensitivity} summarizes the five
main parameters, their possible changes, the projected effect on the
system, and the observations that support each assessment.

\begin{table}[h]
\centering
\caption{Sensitivity Analysis --- Key System Parameters}
\label{tab:sensitivity}
\renewcommand{\arraystretch}{1.35}
\begin{tabular}{>{\bfseries\small}p{0.15\linewidth}
                >{\small}p{0.17\linewidth}
                >{\small}p{0.23\linewidth}
                >{\small}p{0.30\linewidth}}
\toprule
\textbf{Parameter} & \textbf{Possible Change} &
\textbf{System Effect} & \textbf{Observation} \\
\midrule
Number of reports
  & Sharp increase
  & Verification slows; alert latency rises
  & 35\% of surveyed students reported an incident, suggesting
    high potential report volume \\
\rowcolor{rowgray}
User participation
  & Low engagement
  & Fewer incidents logged; real risks missed
  & Low reporting culture observed across all data methods \\
Time of day
  & Night-time rise
  & Higher incident frequency; more alerts per hour
  & All suspicious events observed during evening or night hours \\
\rowcolor{rowgray}
Location
  & High-risk campus zones
  & Localized spike in alerts for specific areas
  & Parking lot identified as highest-risk zone across all sessions \\
Weather conditions
  & Rain or adverse weather
  & Campus movement changes; visibility drops
  & Both suspicious events coincided with rainy night conditions \\
\bottomrule
\end{tabular}
\end{table}

\subsection{Feedback Loops and Emergent Behavior}

The system exhibits both reinforcing and balancing feedback dynamics.
The \textbf{reinforcing loop} works as follows: when users receive
alerts, they tend to become more vigilant and more likely to report
future suspicious activities, which generates more alerts, which
increases awareness further. As adoption grows, the system becomes
progressively more effective \cite{meadows2008}. The \textbf{balancing
loop} works in the opposite direction: if too many alerts are sent,
users may start ignoring them --- a well-documented phenomenon known
as \textit{alert fatigue} --- which reduces system effectiveness over
time. This is why the verification mechanism is a critical design
element: by filtering low-confidence reports before broadcasting, it
preserves signal quality and keeps the alert stream meaningful. Two
additional non-linear effects are worth noting. First, the relationship
between user adoption and report volume is non-linear: a small increase
in active users can produce a disproportionately large jump in
simultaneous report volume during high-activity periods. Second, spatial
and temporal clustering means that alert activity is not uniformly
distributed across the campus or the day, so the system should apply
zone-weighted and time-weighted priority logic rather than treating all
incoming reports identically \cite{checkland1981}.

% ============================================================
%  VI. DISCUSSION
% ============================================================
\section{Discussion}

Three findings emerged consistently across all data collection methods.
The \textbf{reporting culture gap} is the most actionable: most security
incidents go unreported through official channels. The barrier is not
the absence of events but the friction of current reporting tools. A
mobile-first, low-threshold interface directly targets this problem.
\textbf{Night-time and spatial concentration} means that a uniform
alert policy would be suboptimal. The system should support zone-specific
and time-specific alert configurations to reflect the real risk
distribution observed on campus. The \textbf{verification trade-off}
is the central design tension: the verification step is necessary to
prevent false alerts and alert fatigue, but it introduces latency.
Calibrating the right confirmation threshold --- fast enough for genuine
emergencies, strict enough to suppress noise --- is the most important
design decision for the next phase of the project.

% ============================================================
%  VII. CRITICAL ASSESSMENT
% ============================================================
\section{Critical Assessment}

\subsection{Strengths of the Analysis}

Several aspects of this study stand out as methodologically sound.
The use of six complementary data collection methods --- rather than
relying on a single source --- significantly strengthens the validity
of the findings. The convergence of independent methods (interviews,
surveys, and direct observation) on the same three core problems
(under-reporting, night-time concentration, and absence of a centralized
tool) provides mutual validation that would not be achievable through
any single method alone. Additionally, all data was collected
first-hand by the team, which ensures full traceability and eliminates
intermediary distortion. The boundary definition is explicit and
operationally useful: by clearly separating what the CSAS controls
from what it must interface with, the analysis avoids scope inflation
and provides a realistic baseline for design. The sensitivity analysis
goes beyond a descriptive account of system elements by identifying
which parameters have the strongest influence on system behavior ---
a distinction that is essential for prioritizing design decisions in
subsequent phases.

\subsection{Limitations and Risks}

Despite these strengths, several limitations constrain the reach of
the conclusions. The survey sample of 20 students is too small to
support statistical generalization to the full university population.
While the 35\% incident rate is a meaningful indicator, it cannot be
treated as a representative prevalence estimate without a larger and
more stratified sample. The observation phase covered only four zones
in eight sessions; indoor corridors, cafeterias, and transit zones
between buildings were excluded, leaving potential blind spots in the
risk map. The document analysis was limited by restricted access to
detailed institutional records: the absence of zone- and time-specific
breakdown data prevented a more granular historical validation of the
field observations. The weather correlation, while suggestive, is based
on a one-week window and two co-occurring events, which is insufficient
to establish causality or even robust correlation. Any design decisions
based on this relationship should be treated as hypotheses pending
further validation.

\subsection{Gap Between Current and Target System State}

A central finding of this analysis is the measurable gap between the
current state of campus security and the target state the CSAS is
designed to achieve. Table~\ref{tab:gaps} quantifies the most critical
operational gaps identified through the data collection phase.

\begin{table}[h]
\centering
\caption{Operational gap analysis --- current state vs.\ CSAS target state.}
\label{tab:gaps}
\renewcommand{\arraystretch}{1.3}
\begin{tabular}{>{\bfseries\small}p{0.26\linewidth}
                >{\small}p{0.20\linewidth}
                >{\small}p{0.20\linewidth}
                >{\small}p{0.20\linewidth}}
\toprule
\textbf{Dimension} & \textbf{Current State} &
\textbf{CSAS Target} & \textbf{Gap} \\
\midrule
Incident response time
  & $\sim$11.4 min (avg.)
  & $<$5 min
  & $-$6.4 min \\
\rowcolor{rowgray}
Community app adoption
  & $\sim$12\%
  & $\geq$60\%
  & $-$48 pp \\
Incident record completeness
  & $\sim$41\% of events logged
  & $\geq$85\%
  & $-$44 pp \\
\rowcolor{rowgray}
Alert delivery success rate
  & $\sim$87\% (SMS/call)
  & $\geq$99\%
  & $-$12 pp \\
\bottomrule
\end{tabular}
\end{table}

These gaps are not merely technical; they reflect deeply embedded
behavioral and institutional barriers. Closing the adoption gap from
12\% to 60\%, for instance, requires not only a functional mobile
application but a sustained community engagement strategy and
institutional endorsement --- factors that fall outside the strict
technical boundary of the CSAS but are critical to its real-world
effectiveness.

\subsection{Assumptions and Their Implications}

The analysis rests on several assumptions that must be made explicit.
First, it is assumed that students and staff are willing to report
incidents if the friction of doing so is sufficiently reduced. Survey
data supports this directionally (100\% of those who experienced an
incident would use the app), but the assumption has not been validated
under real operating conditions. Second, the verification mechanism is
assumed to be fast enough not to introduce unacceptable latency in
genuine emergencies. This is a critical assumption whose validity
depends entirely on implementation choices --- thresholds, confirmation
logic, and infrastructure capacity --- that have not yet been specified.
Third, the system assumes reliable mobile connectivity across all
campus zones; the observation phase did not systematically assess
signal coverage, and dead zones could materially affect delivery
guarantees. These assumptions represent the primary technical and
behavioral risks for the next design phase.

% ============================================================
%  VIII. CONCLUSIONS
% ============================================================
\section{Conclusions}

This systems analysis established an empirical foundation for the design
of the Community Security Alert System. The boundary definition clearly
separates what the CSAS is responsible for from its external environment.
The stakeholder analysis identified six groups whose needs and
interactions must be balanced in the design. The data collection plan
combined six complementary methods whose findings consistently converged
on the same three key problems. The element and relationship analysis
decomposed the system into five core components and mapped the
information flows between them. The sensitivity analysis showed that
user participation rate and time-of-day are the parameters with the
strongest influence on system behavior. These findings serve as the
baseline for Workshop No.~2, where the system architecture, data model,
and verification logic will be formally specified.

% ── References ───────────────────────────────────────────────
\begin{thebibliography}{9}

\bibitem{sierra2026}
C.~A. Sierra,
\textit{TeamWork as Computer Engineers --- CS Collaboration Guidelines}.
Bogot\'{a}: Universidad Distrital Francisco Jos\'{e} de Caldas, 2026.

\bibitem{checkland1981}
P.~Checkland,
\textit{Systems Thinking, Systems Practice}.
New York: John Wiley \& Sons, 1981.

\bibitem{meadows2008}
D.~H. Meadows,
\textit{Thinking in Systems: A Primer}.
White River Junction, VT: Chelsea Green Publishing, 2008.

\bibitem{iso15288}
ISO/IEC/IEEE,
\textit{Systems and Software Engineering --- System Life Cycle Processes},
ISO/IEC/IEEE Std.~15288:2015, 2015.

\bibitem{creswell2018}
J.~W. Creswell and J.~D. Creswell,
\textit{Research Design: Qualitative, Quantitative, and Mixed Methods
Approaches}, 5th~ed.
Thousand Oaks, CA: SAGE Publications, 2018.

\bibitem{ley1581}
Congreso de Colombia,
\textit{Ley 1581 de 2012 --- R\'{e}gimen General de Protecci\'{o}n
de Datos Personales}, 2012.

\end{thebibliography}

\end{document}
